\subsection*{発展的な読み物}

\addcontentsline{toc}{subsection}{5 発展的な読み物}

\begin{center}
\small
\par\smallskip\noindent\refstepcounter{table}\label{tab:ch14-06}表 \thetable: 発展的な読み物の整理\par\smallskip
表 \thetable は、発展的な読み物の整理を比較・確認しやすい形で整理したものである。\par\smallskip
\begin{longtable}{p{0.38\linewidth}p{0.57\linewidth}}
\toprule
文献 & 読むと得られること \\
\midrule
\endfirsthead
\toprule
文献 & 読むと得られること \\
\midrule
\endhead
Weinberg, The Psychology of Computer Programming & プログラミングを個人とチームの活動として捉える古典的文献である。技術だけでなく、人の認知、チーム、組織がプログラミングに与える影響を学べる。 \\
\midrule
Kinney and Lange, Intellectual Property Law for Business Lawyers & 米国を中心に、知的財産法がどのような考え方で構成されているかを学べる。ソフトウェア契約、著作権、特許、営業秘密を考える入口になる。 \\
\midrule
\bottomrule
\end{longtable}
\normalsize
\end{center}
